\documentclass[10pt,a4paper]{amsart}
\usepackage[margin=19mm]{geometry}
\usepackage{amsmath,amssymb,mathtools}
\usepackage[scr=rsfs,cal=euler]{mathalfa}
\usepackage{xfrac}
\usepackage[unicode,hidelinks,bookmarks=false]{hyperref}
\newcommand{\ie}{i.e.\ }
\newcommand{\cf}{cf.\ }
\newcommand{\resp}{respectively\ }
\newcommand{\Subsection}{\subsection}
\providecommand{\nobreakdash}{\nobreak\hbox{-}}
\setlength{\emergencystretch}{2em}
\multlinegap0pt
\usepackage{amscd,mathrsfs,txfonts}
\hypersetup{
    colorlinks=true,      
    urlcolor=blue}
\usepackage{tikz,pgfplots,pgf}
\usepackage{mathtools}
\usepackage{tikz-cd}
\usepackage{xy}
\usepackage{lipsum}
\newcommand\blfootnote[1]{%
  \begingroup
  \renewcommand\thefootnote{}\footnote{#1}%
  \addtocounter{footnote}{-1}%
  \endgroup
}
\tikzset{node distance=3cm, auto}
\usetikzlibrary{calc} 
\usetikzlibrary{trees} 
\newtheorem{theorem}{Theorem}[section]
\newtheorem{proposition}[theorem]{Proposition}
\newtheorem{definition}[theorem]{Definition}
\newtheorem{corollary}[theorem]{Corollary}
\newtheorem{problem}{Problem}
\newtheorem{lemma}[theorem]{Lemma}
\newtheorem{remark}[theorem]{Remark}
\newtheorem{remarks}[theorem]{Remarks}
\newtheorem{example}[theorem]{Examples}
\newtheorem{conjecture}[theorem]{Conjecture}
\def\A{\mathcal{A}}
\def\B{\mathcal{B}}
\def\F{\mathcal{F}}
\def\G{\mathcal{G}}
\def\H{\mathcal{H}}
\def\I{\mathcal{I}}
\def\J{\mathcal{J}}
\def\C{\mathcal{C}}
\def\L{\mathcal{L}}
\def\M{\mathcal{M}}
\def\R{\mathbb{R}}
\def\S{\mathcal{S}}
\def\U{\mathcal{U}}
\def\T{\mathcal{T}}
\def\C{\mathcal{C}}
\def\K{\mathbb{K}}
\def\N{\mathbb{N}}
\def\linG{\mathrm{lin}(\Gamma(\mathbb{D}))}
\def\Lip{\mathrm{Lip}_0}
\def\Lipnorm{\mathrm{Lip}}
\def\Cq{\mathcal{C}_q^{\mathrm{Lip}_0}}
\def\Cqnorm{\mathcal{C}_q^{\mathrm{Lip}}}
\def\abco{\mathrm{abco}}
\def\co{\mathrm{co}}
\def\Tp{\mathcal{T}_p^{\mathrm{Lip}_0}}
\def\Tpnorm{\mathcal{T}_p^{\mathrm{Lip}}}
\def\lin{\mathrm{lin}}
\def\Rad{\mathrm{Rad}}
\def\RAD{\mathrm{RAD}}
\allowdisplaybreaks
\begin{document}
\title[Lipschitz stable sequence classes: an approach to Rademacher]{Lipschitz stable sequence classes: an approach to Rademacher type and cotype of Lipschitz functions}
\author{D. Ruiz-Casternado}
\date{}
\maketitle
\noindent\emph{Source and licence.} Lipschitz stable sequence classes: an approach to Rademacher type and cotype of Lipschitz functions. Version arXiv:2606.04633v1 (CC BY 4.0). This material is distributed under the Creative Commons Attribution 4.0 International licence, http://creativecommons.org/licenses/by/4.0/, as recorded in the arXiv OAI licence metadata for this exact version and shown on the abstract page https://arxiv.org/abs/2606.04633v1. Full rights notice: licenses/arxiv-2606.04633-CC-BY-4.0.txt.\par\smallskip
\noindent\textit{Editorial scope.} Counted unit: the original \texttt{theorem} environment \texttt{Th 1} at source lines 178--194 together with its complete proof at lines 196--248; the delivered document keeps source lines 170--249 so that every referenced label and every citation resolves inside it. Native editable source is the paper's own concatenated TeX; the AMS wrapper is editorial formatting only. No publisher class, upstream script or unrelated artwork is redistributed.
\section{Lipschitz stability of sequence classes}\label{Sec 2}

Let us recall by \cite[Definition 2.1]{BotCam-17} that a sequence class $Z$ is an assignment that associates to each Banach space $E$ a vector subspace of $E^\N$, denoted $Z(E)$, endowed with a complete norm $\left\|\cdot\right\|_{Z(E)}$ such that $c_{00}(E) \subseteq Z(E) \subseteq \ell_\infty(E)$ with $\left\|\cdot\right\|_\infty \leq \left\|\cdot\right\|_{Z(E)}$, and $\|(e_n^{(j)})_{n\in\N}\|_{Z(\K)}=1$ for all $j \in \N$. Particular examples of sequence classes can be found in \cite[Example 2.2]{BotCam-17}.

If for any $(x_n)_{n\in\N} \in E^\N$ we have that $(x_n)_{n\in\N}\in Z(E)$ if and only if $\sup_{N\in\N} \|(x_k)_{k=1}^N\|_{Z(E)}<\infty$, then the sequence class $Z$ is called finitely determined. In such a case $\|(x_n)_{n\in\N}\|_{Z(E)}=\sup_{N\in\N}\|(x_k)_{k=1}^N\|_{Z(E)}$. $\RAD,\ell_p$ and $\ell_p^w$ for $p\in [1,\infty]$ are examples of finitely determined sequence classes.

Firstly, we analyse the transformation of Banach-valued sequences by means of Lipschitz functions.


\begin{theorem}\label{Th 1}
	Let $f\in \Lip(X,E)$ and let $Z,Y$ be sequence classes. Consider the following assertions:
	\begin{itemize}
		\item[$(i)$] If $(x_n)_{n\in\N},(y_n)_{n\in\N}\in Z(X)$ then $(f(x_n)-f(y_n))_{n\in\N}\in Y(E)$.
		\item[$(ii)$] The operator $\hat{f}: Z(X)\to Y(E)$ given by
		$$
		\hat{f}((x_n)_{n\in\N}) = (f(x_n))_{n\in\N} \qquad ((x_n)_{n\in\N} \in Z(X))
		$$
		is well-defined.
		\item[$(iii)$] $\hat{f} \in \Lip(Z(X),Y(E))$.
		\item[$(iv)$] There exists $C>0$ so that for all $N\in \N$ and all $(x_1,\ldots,x_N)$, $(y_1,\ldots,y_N)\in X^N$, we have
		$$
		\left\|(f(x_k)-f(y_k))_{k=1}^N\right\|_{Y(E)} \leq C\left\|(x_k-y_k)_{k=1}^N\right\|_{Z(X)}.
		$$
	\end{itemize}
	Then $(i) \Leftrightarrow (ii) \Leftarrow (iii) \Rightarrow (iv)$.  Moreover, if the sequence classes $Z$ and $Y$ are finitely determined, then $(iv) \Rightarrow (iii)$. Moreover, in this case, $\Lipnorm(\hat{f}) = \inf\{C\}$ with the infimum being taken over all such constants as in $(iv)$.
\end{theorem}

\begin{proof}
	$(i) \Rightarrow (ii):$ We begin by showing that $\hat{f}$ is a function. Towards this end, let $(x_n)_{n\in\N},(y_n)_{n\in\N} \in Z(X)$ be such that $\hat{f}((x_n)_{n\in\N}) \neq \hat{f}((y_n)_{n\in\N})$. Then $(f(x_n))_{n\in\N} \neq (f(y_n))_{n\in\N}$, i.e., there is $k\in\N$ for which $f(x_k)\neq f(y_k)$, and since $f$ is a function, we have $x_k \neq y_k$. Hence $(x_n)_{n\in\N} \neq (y_n)_{n\in\N}$. Moreover, we claim that $\hat{f}$ is well-defined. Indeed, let $(x_n)_{n\in\N} \in Z(X)$. Since $Z(X)$ is a vector space, the zero sequence $(0)_{n\in\N} \in Z(X)$. Thus applying our hypothesis we have
	$$
	\hat{f}((x_n)_{n\in\N})=(f(x_n))_{n\in\N}=(f(x_n)-f(0))_{n\in\N} \in Y(E).
	$$
	Consequently, $\hat{f}$ is a well-defined operator.

	$(ii) \Rightarrow (i):$ Assume that $\hat{f}:Z(X) \to Y(E)$ is a well-defined operator. Let $(x_n)_{n\in\N},(y_n)_{n\in\N} \in Z(X)$ and note that
	$$
	(f(x_n)-f(y_n))_{n\in\N} = \hat{f}((x_n)_{n\in\N})-\hat{f}((y_n)_{n\in\N}) \in Y(E).
	$$

	$(iii) \Rightarrow (i):$ Let us suppose that $\hat{f} \in \Lip(Z(X),Y(E))$, and let $(x_n)_{n\in\N},(y_n)_{n\in\N} \in Z(X)$. Then
	$$
	\|(f(x_n)-f(y_n))_{n\in\N}\|_{Y(E)}=\|\hat{f}((x_n)_{n\in\N})-\hat{f}((y_n)_{n\in\N})\|_{Y(E)} \leq \Lipnorm(\hat{f}) \|(x_n)_{n\in\N}-(y_n)_{n\in\N}\|_{Z(X)} < \infty,
	$$
	and we obtain that $(f(x_n)-f(y_n))_{n\in\N} \in Y(E)$.

	$(iii) \Rightarrow (iv):$ Assume that $\hat{f}\in \Lip(Z(X),Y(E))$. Let $N\in\N$ and $(x_1,\ldots,x_N)$, $(y_1,\ldots,y_N)\in X^N$. Let us define the sequences $(u_n)_{n\in\N}$ and $(v_n)_{n\in\N}$ by
	$$
	u_k = \begin{cases}
		x_k \quad &k \leq N,\\
		0 \quad &k > N.
	\end{cases} \qquad
	v_k = \begin{cases}
		y_k \quad &k \leq N,\\
		0 \quad &k > N.
	\end{cases}
	$$

	Since $c_{00}(X) \subseteq Z(X)$ then $(u_n)_{n\in\N},(v_n)_{n\in\N} \in Z(X)$. Therefore,
	$$
	\hat{f}((u_n)_{n\in\N})-\hat{f}((v_n)_{n\in\N}) = (f(u_n)-f(v_n))_{n\in\N} = (f(x_1)-f(y_1), \ldots, f(x_N)-f(y_N), 0, \ldots),
	$$
	and we deduce that
	\begin{align*}
		\|(f(x_k)-f(y_k))_{k=1}^N\|_{Y(E)} &= \|\hat{f}((u_n)_{n\in\N})-\hat{f}((v_n)_{n\in\N})\|_{Y(E)}\\ 
		&\leq \Lipnorm(\hat{f})\|(u_n-v_n)_{n\in\N}\|_{Z(X)}\\ 
		&= \Lipnorm(\hat{f})\|(x_k-y_k)_{k=1}^N\|_{Z(X)}.
	\end{align*}

	Hence we obtain the desired result and $\inf\{C\} \leq \Lipnorm(\hat{f})$, where the infimum is taken over all such positive constants for which the previous inequality is fulfilled.

	$(iv) \Rightarrow (iii):$ Suppose that there exists a constant $C>0$ so that
	$$
	\left\|(f(x_k)-f(y_k))_{k=1}^N\right\|_{Y(E)} \leq C\left\|(x_k-y_k)_{k=1}^N\right\|_{Z(X)}
	$$
	for all $N\in\N$ and all $(x_1,\ldots,x_N)$, $(y_1,\ldots,y_N)\in X^N$. Since $Z$ and $Y$ are finitely determined, just taking the supremum over all such $N\in\N$ we obtain
	$$
	\|\hat{f}((x_n)_{n\in\N})-\hat{f}((y_n)_{n\in\N})\|_{Y(E)} \leq C\|(x_n)_{n\in\N}-(y_n)_{n\in\N}\|_{Z(X)}.
	$$
	Hence $\hat{f} \in \Lip(Z(X),Y(E))$ and $\Lipnorm(\hat{f}) \leq \inf\{C\}$ with the infimum being taken over all such positive constants satisfying that inequality.
\end{proof}


\begin{thebibliography}{99}
\bibitem[BotCam-17]{BotCam-17} G. Botelho, J. R. Campos, On the transformation of vector-valued sequences by linear and multilinear operators, Monatsh Math. \textbf{183} (2017), 415--435. \url{https://doi.org/10.1007/s00605-016-0963-4}
\end{thebibliography}
\end{document}
